% TOP_PROBLEM: {"id": 2819, "title": "Kirby Problem 3.21", "category_id": 7, "category": {"id": 7, "name": "topology", "display_name": "Topology", "description": "Properties preserved under continuous deformations.", "slug": "topology", "order_index": 7, "created_at": "2026-07-31T15:26:25.671Z"}, "status": "open"}
% FIRST_POSTED: null
% ENTRY_KIND: statement_only
% Imported from: batch12.tex; UnsolvedMath ID 2819
\documentclass[11pt]{article}
\usepackage[margin=25mm]{geometry}
\usepackage{fontspec}
\setmainfont{Latin Modern Roman}
\usepackage{amsmath,amssymb,mathtools}
\usepackage{unicode-math}
\setmathfont{Latin Modern Math}
\usepackage{xurl,hyperref}
\hypersetup{colorlinks=true,urlcolor=blue}
\setcounter{secnumdepth}{0}
\setlength{\parindent}{0pt}
\setlength{\parskip}{5pt}
\setlength{\emergencystretch}{3em}
\newcommand{\sref}[2]{\hyperlink{source:#1:#2}{[S#2]}}
\newcommand{\cataloguescope}[1]{\paragraph{Recorded target.}#1}
\begin{document}
% UnsolvedMath ID 2819; source catalogue code KP-3.21
\setcounter{section}{2818}
\section{Kirby Problem 3.21}\label{2819}
{\small\sffamily UnsolvedMath ID 2819 · Topology}
\subsection{Definitions and mathematical statement}

\paragraph{Shared notation.}$\mathbb N=\{1,2,\ldots\}$, and $\mathbb Z,\mathbb Q,\mathbb R,\mathbb C$ denote the integers, rationals, reals and complex numbers. Any other conventions are specified locally.


All manifolds here are orientable. A flow is transitive when some orbit is dense. Two flows on the same oriented manifold are orbit equivalent if an orientation-preserving homeomorphism sends their orbits to one another. An Anosov flow generated by $X$ has an invariant splitting $TM=E^s\oplus\mathbb RX\oplus E^u$ and constants $C,a>0$ giving exponential contraction on $E^s$ forwards and on $E^u$ backwards. A pseudo-Anosov flow allows finitely many singular periodic orbits with at least three prongs in its stable/unstable foliations, while retaining the hyperbolic contraction/expansion elsewhere. A quasigeodesic flow has a uniform quasigeodesic bound for its arclength-parametrized lifted orbits in the universal cover, with respect to the lift of any fixed Riemannian metric on the compact manifold. A perfect fit is the stable/unstable missing-corner configuration in the lifted flow orbit space: a stable and an unstable leaf are disjoint but form adjacent sides of a product rectangle with their common corner removed.
For triangulations use the transverse veering convention in the cited Agol--Tsang construction. A taut ideal triangulation has angles $0$ or $\pi$, two opposite $\pi$ edges per tetrahedron, and angle sum $2\pi$ at each edge. A transverse taut structure also coorients its faces compatibly under gluings, with exactly one of the two faces incident to any $0$-angle edge pointing out of its tetrahedron. Color edges red or blue; viewing each flattened tetrahedron from above according to these coorientations, its four outer edges are red, blue, red, blue counterclockwise, starting at the front endpoint of its inner edge. There is no color restriction on the top and bottom $\pi$ edges. Triangulations, including these structures, are counted up to combinatorial equivalence. \sref{2819}{3} A cusped finite-volume hyperbolic manifold is the interior of a compact manifold with torus boundary and a complete finite-volume curvature-$-1$ metric.
\paragraph{Statement recorded in the supplied source.}
For every closed orientable $3$-manifold $M$, are each of the following sets of transitive flows finite up to orbit equivalence: (a) Anosov flows; (b) pseudo-Anosov flows; (c) quasigeodesic pseudo-Anosov flows; (d) pseudo-Anosov flows without perfect fits? Independently, (e) does every orientable finite-volume cusped hyperbolic $3$-manifold have only finitely many veering triangulations? \sref{2819}{2}
\subsection{Short English statement}
Are there finitely many transitive flows of each of four specified types on any fixed closed orientable three-manifold, and finitely many veering triangulations on any fixed cusped hyperbolic three-manifold?
\subsection{Sources}
\begin{description}
\item[\hypertarget{source:2819:1}{S1}] ULAM.AI and the UnsolvedMath Contributors, \emph{UnsolvedMath: A Curated Collection of Open Mathematics Problems}, record 2819 (KP-3.21). CC BY 4.0. \url{https://huggingface.co/datasets/ulamai/UnsolvedMath}.
\item[\hypertarget{source:2819:2}{S2}] R. İnanç Baykur, Robion C. Kirby and Daniel Ruberman (editors), \emph{K3: A New Problem List in Low-Dimensional Topology}, Mathematical Surveys and Monographs 295 (2026), Problem 3.21, p. 145 and following remarks; original author version. \url{https://math.berkeley.edu/sites/default/files/surv-295-ruberman-watermarked-author-pdf.pdf}.
\item[\hypertarget{source:2819:3}{S3}] Ian Agol and Chi Cheuk Tsang, \emph{Dynamics of veering triangulations: infinitesimal components of their flow graphs and applications}, Algebraic and Geometric Topology 24 (2024), 3401--3453; arXiv:2201.02706v2, Definitions 2.1--2.2 and Remark 2.3. \url{https://arxiv.org/abs/2201.02706v2}.
\end{description}
\label{end:2819}
\end{document}
